Fix an integer $n\ge 1$. Let $P$ be a set of size
\[
|P| = n^2+n+1.
\]
A \emph{finite projective plane of order $n$} (as defined in the problem statement) is a collection $\mathcal L$ of subsets of $P$ (called \emph{lines}) such that:

(a) Every line has size $n+1$: for all $L\in\mathcal L$, $|L|=n+1$.

(b) Every unordered pair of distinct points lies on exactly one line: for all distinct $x,y\in P$, there exists a unique $L\in\mathcal L$ with $\{x,y\}\subseteq L$.

Question: If such $(P,\mathcal L)$ exists, must $n$ be a \emph{prime power} (i.e. $n=p^m$ for some prime $p$ and integer $m\ge 1$)?

Edge cases/conventions: The definition makes sense for $n=1$ and gives a non-degenerate structure. We do not consider $n=0$.
